% Independent audit fragment for the corrected array-truth interface.
% Reviewed files: kernel_notes/array_truth_checker{,_bytes,_joint}.py.
% IMPORTANT: the final mode comparison is with u, not s.
% Requires only ordinary LaTeX plus amsmath/amssymb.

\section{Why the absent major can be a proved empty array}
\label{sec:neutral-major-audit}

This section audits the new control-flow change on raw input, including
noncanonical JSON.  It uses the inherited sorted arity decoder and
semantic-opening lemmas for the numerical and logical rules.  It does not
replace those lemmas with finite testing or claim proof-assistant
certification.

\subsection{The exact change}

The initial context is the one-element list \texttt{[[]]}.  Its sole
formula, the empty array, denotes $0=0$ under the inherited decoder.
At every non-hypothesis command the checker now executes
\begin{verbatim}
    m, s, u, t = q()
    c(m, g)
\end{verbatim}
without testing whether \texttt{m} is nonempty.  A canonical command with
no substantive major puts \texttt{[]} in this field.  Its premise closes
immediately by hypothesis, without consuming a command or changing the
fresh-name floor.

The final choice of the retained assumption uses
\begin{verbatim}
    (m or len(a) % 2 == u) and k
\end{verbatim}
instead of comparing the arity parity with \texttt{m}.  In an
introduction command, \texttt{u=0} selects equation progress and
\texttt{u=[]} selects the zero-case premise.  This reuse is legitimate
because introduction is reached only when both \texttt{m} and
\texttt{u} are false-valued.  Computation still requires a truthy
initializer \texttt{u}; its missing literal-zero cases have the
inherited certified-zero repair.

The source is deliberately left with the old matcher
\begin{verbatim}
    m in ([s, u], not s and [t] * 3)
\end{verbatim}
and no added \texttt{m and} guard.  As an isolated Python expression this
would be unsafe: if \texttt{s} is truthy, its second alternative is
\texttt{False}, which compares equal to numeric zero.  The following
whole-program invariant proves that this dangerous major cannot reach
the matcher.

\subsection{Two facts about the root opener}

Call a Python value \emph{numeric} here when it is an integer, a float,
or a Boolean returned by the JSON loader.  The finite Python execution
also admits the loader's optional nonfinite float values.  No custom
numeric class is supplied by a certificate.

Write $O_k(r;t)$ for the opener on the numeric value $r$, with image $t$
and its default formula-root depth zero.  Direct evaluation of the
program gives
\[
 O_k(r;t)=
 \begin{cases}
 r,&k\le r<0,\\
 0,&\text{otherwise}.
 \end{cases}
 \tag{Root scalar}
\]
In particular this result is numeric and independent of the image $t$.
The formula covers negative allocated names and negative fractional
names without identifying different numerical values.  Signed zero and
Booleans have Python's ordinary numerical equality behavior.  A NaN
fails the interval test and returns zero; the relevant infinite values
also return zero.  None of these results is a list.

A second fact is that a successful root opening of a list returns a
list.  Under either strict list predicate, no other input type can
return a numeric result except a numeric input.  Under the generic
iterable opener, a successful iterable traversal also returns a list;
if traversal fails, the numerical comparison cannot succeed on a list,
string, dictionary, or null.  Thus in every profile:
\begin{quote}
If a root opening returns a numeric result, its argument was numeric;
if it succeeds on a list, its result is a list.
\end{quote}
These statements concern finite successful evaluations.  They assert
neither arbitrary revalidation idempotence at ignored annotations nor
unlimited-stack termination for the generic iterable profile.

\subsection{No numeric goal can be proved from list-valued assumptions}

\paragraph{Lemma.}
For any initial integer fresh-name floor and any finite remaining
certificate, there is no finite successful invocation
\texttt{c(a,g)} whose raw goal \texttt{a} is numeric and whose every
context element is a list.

\emph{Proof.}
Suppose otherwise, and select such a successful invocation with the
fewest calls to \texttt{c} in its completed recursive trace.  Its first
hypothesis check cannot close: the root opening of its numeric goal is
numeric, and no numerical JSON value equals a Python list.

The invocation must therefore consume a command and successfully return
from \texttt{c(m,g)}.  If \texttt{m} were numeric, this would be a
strictly shorter successful invocation of the forbidden kind.  Thus the
major is not numeric.

In either strict profile, a successfully checked nonnumeric major is a
list.  In the generic profile an empty string or an empty dictionary can
also open to the empty list.  A nonempty dictionary can occasionally
have its iterable alias proved, but it cannot occur as a truthy major
in the selected completed parent invocation: the subsequent expression
\texttt{m * 3} raises a type error.  A nonempty string has no finite
successful opening.  Consequently every truthy major that survives to
the continuation exports a list, and every surviving false-valued
major is an empty list or one of those harmless empty iterable aliases.

Consider the possible change of goal.  If the matcher does not fire,
the goal stays numeric.  If it fires, its target test requires
\texttt{i(t,u) == a}.  The second root-opening fact therefore implies
that \texttt{t} is numeric.  By (Root scalar), \texttt{i(t,s)} and
\texttt{i(t,u)} are the same root result.  Both the Boolean-expression
and the explicit-conditional implementations consequently retain a
numeric goal.  The Boolean implementation may retain the original raw
zero alias when the transported result is false-valued; that makes no
difference to this assertion.

If the major is truthy, the generated hypothesis is the list
\texttt{m * 3}.  If the major is false-valued and \texttt{u} is truthy,
the computation expression constructs a two-element list.  A successful
opening of either hypothesis is a list.  The final recursive call then
has a numeric root-opened goal and an entirely list-valued context.  It
is a proper subtrace of the selected invocation, contradicting its
minimality.

The only remaining case is introduction.  Its expression
\texttt{i(*a,t,2)} must unpack the numeric raw goal as an iterable,
which raises a type error.  That case cannot belong to a successful
trace.  All possibilities have been excluded.\hfill$\square$

The minimization ranges over all starting floors and stacks.  It is
therefore unaffected by a recursively checked major decreasing the
global floor before its parent resumes.  The argument uses finite
completed traces, not a presumption that recursive checking always
terminates.

\subsection{Every reached command has a list-valued context}

Initially the only assumption is the list \texttt{[]}.  Suppose a
command is reached with list-valued assumptions and its major check
returns.  The lemma excludes a numeric major.  A truthy surviving major
exports a list; a computation constructs a list.  In introduction the
assignment is
\begin{verbatim}
    h = a = i(*a, t, 2)
\end{verbatim}
in mathematical order.  If its result is numeric, the later
\texttt{len(a)} raises a type error before the continuation context is
constructed, since this branch has false-valued \texttt{m}.  Such a
scalar matrix does \emph{not} insert a scalar-zero assumption.  If the
assignment yields a list, successful opening again yields a list.

Both recursive premises thus receive list-valued contexts whenever
execution reaches them.  The expression \texttt{g + [new]} constructs a
new list; a premise cannot mutate its caller's context.  Induction on
the reached calls proves the invariant.  The tuple-based explanatory
variant has the corresponding immutability property.

It follows that no numeric zero, Boolean \texttt{False}, floating zero,
truthy numeric scalar, or newly allocated numeric name can be a
successfully verified major at a reached command.  In particular a
closed scalar-zero goal is rejected by this interface.  The canonical
empty-array truth is accepted instead; mathematical completeness is
claimed for canonical encodings, not for every redundant truth alias.

\subsection{The unguarded matcher is safe on every reached continuation}

If \texttt{s} is false-valued, both matcher alternatives are nonempty
lists, so no false-valued major can match.  If \texttt{s} is truthy, the
second alternative is \texttt{False}.  A false-valued major reaching
this point is an empty list, or in the generic profile an empty string
or dictionary.  None equals \texttt{False}; a numeric major has already
been excluded.  Therefore, at every point reached after a successful
major check, the matcher has exactly the same truth value as
\begin{verbatim}
    m and m in ([s, u], not s and [t] * 3)
\end{verbatim}
even though this identity is false on unrestricted Python tuples of
values.  No sentinel and no extra source guard are needed.

The proof depends on the initial context being \texttt{[[]]}, not
\texttt{[0]}, and on premise contexts being freshly constructed.
Silently changing either detail would invalidate this argument.

\subsection{Why the relocated induction mode must be the initializer field}

The corrected selector compares arity parity with \texttt{u}.  In an
introduction branch this field is necessarily false-valued.  No
false-valued JSON value compares equal to one.  Thus a meaningful
compound matrix, of odd arity one or three, cannot receive a progress
assumption.  A literal equation has even arity two and receives
progress precisely for a zero-like mode such as \texttt{0}; an empty
array mode chooses the zero instance.  Surplus arities retain their
inherited constant interpretations.  Scalar matrices cannot survive
the shape check discussed above.

Comparing instead with the unused \texttt{s} field would be wrong.
That field is unrestricted in introduction, so choosing
\texttt{s=1} would enable progress for odd-arity compound matrices.
The distinction is semantic, not merely a preferred serialization.

In the computation branch \texttt{u} is truthy.  Its new comparison can
choose a fresh-name opening for some noncanonical scalar initializer
aliases.  This is handled by the inherited zero-instance discharge
lemma, not by claiming that every generated open hypothesis is
uniformly true.  If the generated hypothesis is $H(\vec b)$, its zero
completion $H(\vec0)$ is a computation theorem, while the old-floor
validation makes the finite family of newly admitted names $\vec b$
fresh for the retained conclusion and context.  Substituting zero for
those names and cutting in the computation theorem discharges the
extra hypothesis.  Negative fractional names in the newly admitted
interval are included separately in this finite substitution.

\subsection{Conservativity and canonical completeness}

The initial empty-array formula is already a theorem.  For example the
permitted nonempty-seed computation gives
\[
 C=0,\qquad C:=I(S(S0),S0;v.0).
\]
Use this equation as both the equality major and the premise of
transport through $F(w)\equiv(w=0)$ to obtain $0=0$.  Adding that fixed
theorem to a proof context is conservative by substitution of its
finite derivation for hypothesis uses.  No new object-level truth
predicate or logical axiom has been introduced.

For canonical certificates, substantive major commands retain their
old fields.  A computation uses the absent-major field \texttt{[]} and
its existing nonempty initializer.  An introduction uses absent major
\texttt{[]}, with \texttt{u=0} for equation progress and
\texttt{u=[]} for cases; its witness remains in \texttt{t}.
The added empty major closes without a command or name allocation, so
the remaining canonical proof simulation is unchanged.  First applying
the existing certified-zero repair supplies every unrestricted
computation instance.  This gives a finite certificate for every
canonical derivation of the inherited native calculus.

Conversely, the invariant makes each reached matcher a legitimate
guarded matcher.  The remaining commands use the inherited
transport, closure, paired-introduction/elimination, and zero-instance
computation arguments, with the corrected equation-only mode test.
Discharging the initial $0=0$ theorem yields a derivation in that same
calculus.  Its established whole-domain iPRA interpretations therefore
transfer without changing the numerical domain, equality, or object
language.  The new result concerns a proof-checker presentation, not a
new primitive-recursive arithmetic foundation.

\subsection{Raw JSON types and the generic profile}

The strict profiles recognize only Python lists as containers.  The
short predicate \texttt{a * 0 == []} has this effect for JSON values:
lists pass; strings and numbers give different zero products;
dictionaries and null raise.  Its numerical fallback rejects strings.
The root invariant covers all numbers and Booleans that reach it,
including nonfinite values accepted by the default Python loader.

The generic profile has the inherited additional aliases.  An empty
string and an empty dictionary open to \texttt{[]}; a dictionary with
only the empty-string key opens to \texttt{[[]]}, with its value ignored.
A truthy dictionary cannot survive export as a major because
\texttt{m * 3} raises.  Nonempty strings and nonempty dictionary keys
have no finite successful opening.  Actual CPython rejects them through
its recursion limit and subsequent type failure, while an ideal
unlimited-stack execution can diverge.  The raw soundness argument
above covers every finite successful run of this profile; it does not
assert ideal termination on every malformed finite JSON input.

Larger arrays remain arrays throughout opening, irrespective of their
ignored annotations or depth increments.  They therefore do not
threaten the scalar-versus-list invariant.  Their formula meanings and
freshness behavior remain subject to the inherited sorted-decoder
proof.  Physical time, memory, and stack limits can reject sufficiently
large genuine proofs in every profile; the completeness claim has the
same adequate-resource, canonical-proof interpretation as its
predecessors.
